\section{From an integrable section to a rational boundary}
\label{sec:approximation}

The analytic information in Theorem~\ref{thm:controlled-boundary}
concerns one section, rather than a nonvanishing local frame.  The
distinction already appears on a disk.  For a weight
$\phi=a\log|z|^2$ with $a\ge0$ and a section $s_J=z^j$ with
$j\in\mathbb Z_{\ge0}$, the two integrability thresholds are
\[
 |s_J|^2e^{-\phi}\in L^1_{\mathrm{loc}}
       \quad\Longleftrightarrow\quad a<j+1,
 \qquad
 e^{-\phi}\in L^1_{\mathrm{loc}}
       \quad\Longleftrightarrow\quad a<1.
\]
On a resolution, the Jacobian contributes an additional vanishing
order.  After dividing the fixed divisor by the multiple used for
approximation, subtracting this Jacobian order and that of $J$ will
leave coefficients strictly below one.  The resulting signed divisor
has a positive part that can enter a klt boundary, while its negative
part records the
support on which a correction is permitted.

We use additive notation for rational line bundles.  If a metric has
local squared frame norm $e^{-\phi}$, its curvature is positive when
$\phi$ is plurisubharmonic.  A global strict lower bound means
$\sqrt{-1}\partial\bar\partial\phi\ge c\omega$ for one K\"ahler form
$\omega$ and one $c>0$.  Smooth changes of weight do not affect local
integrability.  The multiplier ideal $\mathcal J(\phi)$ consists of
holomorphic germs $u$ for which $|u|^2e^{-\phi}$ is locally integrable.
Section~\ref{sec:correction-ring} uses the unweighted case $J=0$;
it first obtains a global strict bound while preserving integrability
of squared frame norms.

\begin{lemma}[Approximation retaining a distinguished section]
\label{lem:weighted-resolution}
Let $Y$ be smooth projective, let $J\ge0$ be an integral divisor, and
let $B$ be a rational line bundle with a singular metric $h$ having a
global strict curvature lower bound.  Suppose
$\mathcal J(h)\supset\mathcal O_Y(-J)$.
For every sufficiently large divisible integer $r$, the complete system
$|rB|$ is nonempty.  Let $\sigma:Y'\to Y$ be a log resolution of its
base ideal, $J$, and the exceptional locus, and write
\[
 \sigma^*|rB|=G+|\mathrm{Mov}|,
\]
where $G$ is the fixed divisor and $|\mathrm{Mov}|$ is base-point free.
Then the simple-normal-crossing divisor
\begin{equation}\label{eq:weighted-signed-divisor}
 \Psi=G/r-\sigma^*J-K_{Y'/Y}
\end{equation}
has every coefficient strictly less than one.  Its negative part is
supported on $\sigma^*J$ and the exceptional primes.
\end{lemma}

\begin{proof}
Fix a very ample integral divisor $H$ and put $d=\dim Y$.  For $r$
sufficiently large and divisible, the curvature of $rh$ dominates the
fixed smooth curvatures of $K_Y+iH$, simultaneously for $1\le i\le d$.
Nadel vanishing~\cite{Nadel1990,Demailly2015} gives
\[
 H^i\bigl(Y,\mathcal O_Y(rB-iH)\otimes\mathcal J(r\phi)\bigr)=0
 \quad (1\le i\le d).
\]
Indeed, in adjoint form the coefficient line is $rB-K_Y-iH$, whose
metric has a strict positive lower bound.  Castelnuovo--Mumford
regularity makes
$\mathcal O_Y(rB)\otimes\mathcal J(r\phi)$ globally generated.
This is a nonzero coherent sheaf of rank one.  Consequently $|rB|$ is
nonempty and its base ideal $\mathfrak b_r$ satisfies
\begin{equation}\label{eq:base-ideal-multiplier}
 \mathfrak b_r\supset\mathcal J(r\phi).
\end{equation}

We next compare the fixed divisor with the singularity of $h$.  Fix $r$
and a prime $S$ on $Y'$.  At its general point choose a transverse
coordinate $t$ and write
\[
 a=\ord_SG,\qquad j=\ord_S\sigma^*J,
 \qquad k=\ord_SK_{Y'/Y}.
\]
On nested coordinate balls downstairs, the point-extension estimate
of Ohsawa--Takegoshi, in its weighted Bergman formulation
\cite{OhsawaTakegoshi1987,Demailly2015}, gives holomorphic functions
$u_z$ with
\[
 |u_z(z)|^2=e^{r\phi(z)},\qquad
 \int |u_z|^2e^{-r\phi}\le C_r
\]
for every center $z$ in the smaller ball where $\phi(z)>-\infty$.
The constant is independent of $z$; it need not be independent of $r$.
Since $\phi$ is bounded above on a slightly smaller ball, holomorphic
mean-value estimates bound $u_z$ uniformly there.  Each germ of $u_z$
belongs to $\mathcal J(r\phi)$, hence to $\mathfrak b_r$ by
\eqref{eq:base-ideal-multiplier}.  Its pullback is therefore divisible
by $t^a$ near a general point of $S$.  Cauchy estimates on nested
polydisks bound $(u_z\circ\sigma)/t^a$ uniformly as well.  Substituting
$z=\sigma(w)$ proves
\begin{equation}\label{eq:weighted-local-order}
 \phi\circ\sigma\le \frac ar\log|t|^2+O(1).
\end{equation}
The assertion is automatic on the polar set of $\phi$.

The density $|s_J|^2e^{-\phi}$ is locally integrable downstairs.
Change of variables contributes a factor of order $2k$ along $S$.
By \eqref{eq:weighted-local-order}, its pullback is bounded below,
up to a positive bounded factor, by
$|t|^{2(j+k-a/r)}$ at a general point of $S$.  Integrability in the
transverse complex variable forces
\[
 \frac ar-j-k<1.
\]
These are exactly the coefficients in
\eqref{eq:weighted-signed-divisor}.  Resolving the three supports makes
that divisor simple normal crossings.  Away from $\sigma^*J$ and the
exceptional locus, its only possible contribution is $G/r\ge0$,
which proves the assertion about its negative part.
\end{proof}

The strict inequalities are the essential output.  A non-strict
inequality would not produce a klt boundary.  The local extension
argument explains why integrability of the distinguished section is
precisely the hypothesis that survives after subtracting $\sigma^*J$.

\begin{theorem}[Weighted anticanonical approximation]
\label{thm:weighted-approximation}
Let $Y$ be smooth projective and let $J\ge0$ be integral.  Suppose
$-K_Y+J$ has a singular metric with a global strict curvature lower
bound and multiplier ideal containing $\mathcal O_Y(-J)$.
Choose a sufficiently small ample rational divisor $A$.
For a sufficiently large divisible integer $r$, resolve the complete
system $|r(-K_Y+J-A)|$ and $J$ by $\sigma:Y'\to Y$, and write
\[
 \sigma^*|r(-K_Y+J-A)|=G+|\mathrm{Mov}|.
\]
Then $\Psi=G/r-\sigma^*J-K_{Y'/Y}$ has coefficients less than one,
$\Psi^+$ is an effective klt boundary, and
$D_0=\Psi^-$ is effective and supported on $\sigma^*J$ and exceptional
primes.  Moreover
\begin{equation}\label{eq:weighted-adjoint-identity}
 D_0\sim_{\mathbb Q}
 K_{Y'}+\Psi^++\mathrm{Mov}/r+\sigma^*A.
\end{equation}
The rational class $\mathrm{Mov}/r+\sigma^*A$ is nef and big.
If an algebraic torus acts on $Y$, $J$ is invariant, and the indicated
full systems are torus-stable, the resolution may be chosen equivariantly.  Then $G$,
$K_{Y'/Y}$, $\Psi^\pm$ and $D_0$ are invariant, and the free
system $|\mathrm{Mov}|$ is torus-stable.  Identity
\eqref{eq:weighted-adjoint-identity} is an identity of ordinary rational
line bundles; it does not assert that a chosen member of either the
free system or the ample class is invariant.
\end{theorem}

\begin{proof}
Subtracting a sufficiently small smooth ample metric preserves the
strict lower bound and does not change the multiplier ideal.
Apply Lemma~\ref{lem:weighted-resolution} to $B=-K_Y+J-A$.
Its identity $G+\mathrm{Mov}\sim_{\mathbb Q}r\sigma^*B$ gives
\eqref{eq:weighted-adjoint-identity}.  The positive part of a simple
normal crossing divisor with coefficients less than one is klt.
The moving class is nef, and the pullback of the ample class $A$ is
nef and big, proving the last positivity assertion.
In the equivariant setting the complete-system base ideal is invariant;
equivariant resolution in characteristic zero, applied also to $J$,
therefore supplies invariant fixed, relative canonical and boundary
divisors.  This does not require the original metric to be invariant.
\end{proof}

The signed identity supplies more information than an effectivity
conclusion.  In particular, its negative part has an explicit support.
It does not assert integrability of the norm of a nonvanishing frame;
that stronger hypothesis enters Section~\ref{sec:correction-ring}.
